\documentclass[10pt,letterpaper]{article}
\usepackage[utf8]{inputenc}
\usepackage[intlimits]{amsmath}
\usepackage{amsfonts}
\usepackage{amssymb}
\usepackage{ragged2e}
\usepackage[letterpaper, margin=0.5in]{geometry}
\usepackage{graphicx}
\usepackage{cancel}
\usepackage{mathtools}
\usepackage{tabularx}
\usepackage{arydshln}
\usepackage{tensor}
\usepackage{array}
\usepackage{xcolor}
\usepackage[boxed]{algorithm}
\usepackage[noend]{algpseudocode}
\usepackage{listings}
\usepackage{textcomp}
% \usepackage[pdf,tmpdir,singlefile]{graphviz}
\usepackage{mathrsfs}
\usepackage{bbm}
\usepackage{tikz}
\usepackage{tikz-cd}
\usepackage{enumitem}
\usepackage{arydshln}
\usepackage{relsize}
\usepackage{multicol}
\usepackage{scalerel}
\usepackage{upgreek}
\usepackage{wasysym}

\usetikzlibrary{bayesnet}
\setlist{noitemsep}

%%%%%%%%%%%%%%%%%%%%%%%%%%%%%
% Formatting commands
%%%%%%%%%%%%%%%%%%%%%%%%%%%%%
\newcommand{\n}{\hfill\break}
\newcommand{\nn}{\vspace{0.5\baselineskip}\n}
\newcommand{\up}{\vspace{-\baselineskip}}
\newcommand{\hangblock}[2]{\par\noindent\settowidth{\hangindent}{\textbf{#1: }}\textbf{#1: }\nolinebreak#2}
\newcommand{\lemma}[1]{\hangblock{Lemma}{#1}}
\newcommand{\defn}[1]{\hangblock{Defn}{#1}}
\newcommand{\thm}[1]{\hangblock{Thm}{#1}}
\newcommand{\cor}[1]{\hangblock{Cor}{#1}}
\newcommand{\prop}[1]{\hangblock{Prop}{#1}}
\newcommand{\ex}[1]{\hangblock{Ex}{#1}}
\newcommand{\exer}[1]{\hangblock{Exer}{#1}}
\newcommand{\fact}[1]{\hangblock{Fact}{#1}}
\newcommand{\remark}[1]{\hangblock{Remark}{#1}}
\newcommand{\proven}{\;$\square$\n}
\newcommand{\problem}[1]{\par\noindent{#1}\n}
\newcommand{\problempart}[2]{\par\noindent\indent{}\settowidth{\hangindent}{\textbf{(#1)} \indent{}}\textbf{(#1) }\nolinebreak#2\n}
\newcommand{\ptxt}[1]{\textrm{\textnormal{#1}}}
\newcommand{\inlineeq}[1]{\centerline{$\displaystyle #1$}}
\newcommand{\pageline}{\up\par\noindent\rule{\textwidth}{0.1pt}}

%%%%%%%%%%%%%%%%%%%%%%%%%%%%%
% Math commands
%%%%%%%%%%%%%%%%%%%%%%%%%%%%%
% Set Theory
\newcommand{\card}[1]{\left|#1\right|}
\newcommand{\set}[1]{\left\{#1\right\}}
\newcommand{\setmid}{\;\middle|\;}
\newcommand{\ps}[1]{\mathcal{P}\left(#1\right)}
\newcommand{\pfinite}[1]{\mathcal{P}^{\ptxt{finite}}\left(#1\right)}
\newcommand{\naturals}{\mathbb{N}}
\newcommand{\N}{\naturals}
\newcommand{\integers}{\mathbb{Z}}
\newcommand{\Z}{\integers}
\newcommand{\rationals}{\mathbb{Q}}
\newcommand{\Q}{\rationals}
\newcommand{\reals}{\mathbb{R}}
\newcommand{\R}{\reals}
\newcommand{\complex}{\mathbb{C}}
\newcommand{\C}{\complex}
\newcommand{\halfPlane}{\mathbb{H}}
\let\HSym\H
\let\H\relax
\newcommand{\H}{\halfPlane}
\newcommand{\comp}{^{\complement}}
\DeclareMathOperator{\Hom}{Hom}
\newcommand{\Ind}{\mathbbm{1}}
\newcommand{\cut}{\setminus}
\DeclareMathOperator{\elem}{elem}

% Differentiable Manifolds
\newcommand{\RP}{\mathbb{RP}}
\newcommand{\CP}{\mathbb{RP}}
\newcommand{\osubset}{\overset{\mathclap{\scalebox{0.5}{\ptxt{open}}}}{\subset}}
\newcommand{\osubseteq}{\overset{\mathclap{\scalebox{0.5}{\ptxt{open}}}}{\subseteq}}
\newcommand{\osupset}{\overset{\mathclap{\scalebox{0.5}{\ptxt{open}}}}{\supset}}
\newcommand{\osupseteq}{\overset{\mathclap{\scalebox{0.5}{\ptxt{open}}}}{\supseteq}}
\newcommand{\pdat}[3]{\left.\pd{#1}{#2}\right|_{#3}}
\DeclareMathOperator{\so}{so}
\DeclareMathOperator{\codim}{codim}
\DeclareMathOperator{\Diff}{Diff}
\let\dSym\d
\let\d\relax
\newcommand{\d}{\partial}
\DeclareMathOperator{\gl}{gl}

% Graph Theory
\let\deg\relax
\DeclareMathOperator{\deg}{deg}
\newcommand{\degp}{\ptxt{deg}^{+}}
\newcommand{\degn}{\ptxt{deg}^{-}}
\newcommand{\precdot}{\mathrel{\ooalign{$\prec$\cr\hidewidth\hbox{$\cdot\mkern0.5mu$}\cr}}}
\newcommand{\succdot}{\mathrel{\ooalign{$\cdot\mkern0.5mu$\cr\hidewidth\hbox{$\succ$}\cr\phantom{$\succ$}}}}
\DeclareMathOperator{\cl}{cl}
\DeclareMathOperator{\affdim}{affdim}

% Probability
\newcommand{\parSymbol}{\P}
\newcommand{\Prob}{\mathbb{P}}
\renewcommand{\P}{\Prob}
\newcommand{\Avg}{\mathbb{E}}
\newcommand{\E}{\Avg}
\DeclareMathOperator{\Var}{Var}
\DeclareMathOperator{\cov}{cov}
\DeclareMathOperator{\Unif}{Unif}
\DeclareMathOperator{\Binom}{Binom}
\newcommand{\CI}{\mathrel{\text{\scalebox{1.07}{$\perp\mkern-10mu\perp$}}}}

% Standard Math
\newcommand{\inv}{^{-1}}
\newcommand{\abs}[1]{\left|#1\right|}
\newcommand{\ceil}[1]{\left\lceil{}#1\right\rceil{}}
\newcommand{\floor}[1]{\left\lfloor{}#1\right\rfloor{}}
\newcommand{\conj}[1]{\overline{#1}}
\newcommand{\of}{\circ}
\newcommand{\tri}{\triangle}
\newcommand{\inj}{\hookrightarrow}
\newcommand{\surj}{\twoheadrightarrow}
\newcommand{\ndiv}{\nmid}
\renewcommand{\epsilon}{\varepsilon}
\newcommand{\divides}{\mid}
\newcommand{\ndivides}{\nmid}
\DeclareMathOperator{\lcm}{lcm}
\DeclareMathOperator{\sgn}{sgn}
\newcommand{\map}[4]{\!\!\!\begin{array}[t]{rcl}#1 & \!\!\!\!\to & \!\!\!\!#2\\ {}#3 & \!\!\!\!\mapsto & \!\!\!\!#4\end{array}}
\newcommand{\bigsum}[2]{\smashoperator[lr]{\sum_{\scalebox{#1}{$#2$}}}}
\DeclareMathOperator{\gcf}{gcf}
\newcommand{\restr}[2]{\left.#1\right|_{#2}}

% Linear Algebra
\newcommand{\Id}{\textrm{\textnormal{Id}}}
\newcommand{\im}{\textrm{\textnormal{im}}}
\newcommand{\norm}[1]{\abs{\abs{#1}}}
\newcommand{\tpose}{^{T}\!}
\newcommand{\iprod}[1]{\left<#1\right>}
\newcommand{\giprod}{\iprod{\;\,,\;}}
\DeclareMathOperator{\tr}{tr}
\DeclareMathOperator{\trace}{tr}
\newcommand{\chgBasMat}[3]{\!\!\tensor*[_{#1}]{\left[#2\right]}{_{#3}}}
\newcommand{\vecBas}[2]{\tensor*[]{\left[#1\right]}{_{#2}}}
\DeclareMathOperator{\GL}{GL}
\DeclareMathOperator{\Mat}{Mat}
\DeclareMathOperator{\vspan}{span}
\DeclareMathOperator{\rank}{rank}
\newcommand{\V}[1]{\vec{#1}}
\DeclareMathOperator{\proj}{proj}
\DeclareMathOperator{\compProj}{comp}
\DeclareMathOperator{\row}{row}
\newcommand{\smallPMatrix}[1]{\paren{\begin{smallmatrix}#1\end{smallmatrix}}}
\newcommand{\smallBMatrix}[1]{\brack{\begin{smallmatrix}#1\end{smallmatrix}}}
\newcommand{\pmat}[1]{\begin{pmatrix}#1\end{pmatrix}}
\newcommand{\bmat}[1]{\begin{bmatrix}#1\end{bmatrix}}
\newcommand{\dual}{^{*}}
\newcommand{\pinv}{^{\dagger}}
\newcommand{\horizontalMatrixLine}{\ptxt{\rotatebox[origin=c]{-90}{$|$}}}
\DeclareMathOperator{\range}{range}
\DeclareMathOperator{\Symm}{Symm}
\DeclareMathOperator{\SU}{SU}

% Multilinear Algebra
\let\Lsym\L
\let\L\relax
\DeclareMathOperator{\L}{\mathscr{L}}
\DeclareMathOperator{\A}{\mathcal{A}}
\DeclareMathOperator{\Alt}{Alt}
\DeclareMathOperator{\Sym}{Sym}
\newcommand{\ot}{\otimes}
\newcommand{\ox}{\otimes}
\DeclareMathOperator{\asc}{asc}
\DeclareMathOperator{\asSet}{set}
\DeclareMathOperator{\sort}{sort}
\DeclareMathOperator{\ringA}{\mathring{A}}
\DeclareMathOperator{\Sh}{Sh}

% Topology
\newcommand{\closure}[1]{\overline{#1}}
\newcommand{\uball}{\mathcal{U}}
\DeclareMathOperator{\Int}{Int}
\DeclareMathOperator{\Ext}{Ext}
\DeclareMathOperator{\Bd}{Bd}
\DeclareMathOperator{\rInt}{rInt}
\DeclareMathOperator{\ch}{ch}
\DeclareMathOperator{\ah}{ah}
\newcommand{\LargerTau}{\mathlarger{\mathlarger{\mathlarger{\mathlarger{\tau}}}}}
\newcommand{\Tau}{\mathcal{T}}

% Analysis
\DeclareMathOperator{\Graph}{Graph}
\DeclareMathOperator{\epi}{epi}
\DeclareMathOperator{\hypo}{hypo}
\DeclareMathOperator{\supp}{supp}
\newcommand{\lint}[2]{\underset{#1}{\overset{#2}{{\color{black}\underline{{\color{white}\overline{{\color{black}\int}}\color{black}}}}}}}
\newcommand{\uint}[2]{\underset{#1}{\overset{#2}{{\color{white}\underline{{\color{black}\overline{{\color{black}\int}}\color{black}}}}}}}
\newcommand{\alignint}[2]{\underset{#1}{\overset{#2}{{\color{white}\underline{{\color{white}\overline{{\color{black}\int}}\color{black}}}}}}}
\newcommand{\extint}{\ptxt{ext}\int}
\newcommand{\extalignint}[2]{\ptxt{ext}\alignint{#1}{#2}}
\newcommand{\conv}{\ast}
\newcommand{\pd}[2]{\frac{\partial{}#1}{\partial{}#2}}
\newcommand{\del}{\nabla}
\DeclareMathOperator{\grad}{grad}
\DeclareMathOperator{\curl}{curl}
\let\div\relax
\DeclareMathOperator{\div}{div}
\DeclareMathOperator{\vol}{vol}

% Complex Analysis
\let\Re\relax
\DeclareMathOperator{\Re}{Re}
\let\Im\relax
\DeclareMathOperator{\Im}{Im}
\DeclareMathOperator{\Res}{Res}

% Abstract Algebra
\DeclareMathOperator{\ord}{ord}
\newcommand{\generated}[1]{\left<#1\right>}
\newcommand{\cycle}[1]{\smallPMatrix{#1}}
\newcommand{\id}{\ptxt{id}}
\newcommand{\iso}{\cong}
\DeclareMathOperator{\Aut}{Aut}
\DeclareMathOperator{\SL}{SL}
\DeclareMathOperator{\op}{op}
\newcommand{\isom}[4]{\!\!\!\begin{array}[t]{rcl}#1 & \!\!\!\!\overset{\sim}{\to} & \!\!\!\!#2\\ #3 & \!\!\!\!\mapsto & \!\!\!\!#4\end{array}}
\newcommand{\F}{\mathbb{F}}
\newcommand{\acts}{\;\reflectbox{\rotatebox[origin=c]{-90}{$\circlearrowright$}}\;}
\newcommand{\disjunion}{\mathrel{\text{\scalebox{1.07}{$\perp\mkern-10mu\perp$}}}}
\DeclareMathOperator{\SO}{SO}
\DeclareMathOperator{\stab}{stab}
\DeclareMathOperator{\nullity}{nullity}
\DeclareMathOperator{\Perm}{Perm}
\DeclareMathOperator{\nsubgp}{\vartriangleleft}
\DeclareMathOperator{\notnsubgp}{\ntriangleleft}
\newcommand{\presentation}[2]{\left<#1\;\middle|\;#2\right>}
\DeclareMathOperator{\Char}{char}
\DeclareMathOperator{\fchar}{char}
\DeclareMathOperator{\triv}{triv}
\DeclareMathOperator{\reg}{reg}
\DeclareMathOperator{\std}{std}
\DeclareMathOperator{\Func}{Func}
\DeclareMathOperator{\End}{End}

% Convex Optimization
\let\sectionSymbol\S
\let\S\relax
\newcommand{\S}{\mathbb{S}}
\DeclareMathOperator{\dist}{dist}
\DeclareMathOperator{\dom}{dom}
\DeclareMathOperator{\diag}{diag}
\DeclareMathOperator{\ones}{\mathbbm{1}}
\newcommand{\minimizeOver}[3]{\begin{array}{rl}\underset{#1}{\ptxt{minimize}} & #2\\ \ptxt{subject to} & #3\end{array}}
\newcommand{\maximizeOver}[3]{\begin{array}{rl}\underset{#1}{\ptxt{maximize}} & #2\\ \ptxt{subject to} & #3\end{array}}
\newcommand{\minimizationProblem}[2]{\minimizeOver{}{#1}{#2}}
\newcommand{\maximizationProblem}[2]{\maximizeOver{}{#1}{#2}}
\newcommand{\minimizeOverUnconstrained}[2]{\begin{array}{rl}\underset{#1}{\ptxt{minimize}} & #2\end{array}}
\newcommand{\maximizeOverUnconstrained}[2]{\begin{array}{rl}\underset{#1}{\ptxt{maximize}} & #2\end{array}}
\newcommand{\minimizationUnconstrained}[1]{\minimizeOverUnconstrained{}{#1}}
\newcommand{\maximizationUnconstrained}[1]{\maximizeOverUnconstrained{}{#1}}
\DeclareMathOperator{\argmin}{argmin}
\DeclareMathOperator{\argmax}{argmax}

% Proofs
\newcommand{\st}{s.t.}
\newcommand{\unique}{!}
\newcommand{\iffdef}{\overset{\ptxt{def}}{\Leftrightarrow}}
\newcommand{\eqVertical}{\rotatebox[origin=c]{90}{=}}
\newcommand{\mapsfrom}{\mathrel{\reflectbox{\ensuremath{\mapsto}}}}
\newcommand{\mapsdown}{\rotatebox[origin=c]{-90}{$\mapsto$}\mkern2mu}
\newcommand{\mapsup}{\rotatebox[origin=c]{90}{$\mapsto$}\mkern2mu}
\newcommand{\from}{\!\mathrel{\reflectbox{\ensuremath{\to}}}}
\newcommand{\labeledeq}[1]{\overset{\mathclap{\ptxt{#1}}}{=}}
\newcommand{\eqdef}{\labeledeq{def}}

% Brackets
\newcommand{\paren}[1]{\left(#1\right)}
\renewcommand{\brack}[1]{\left[#1\right]}
\renewcommand{\brace}[1]{\left\{#1\right\}}
\newcommand{\ang}[1]{\left<#1\right>}

% Algorithms
\algrenewcommand{\algorithmiccomment}[1]{\hskip 1em \texttt{// #1}}
\algrenewcommand\algorithmicrequire{\textbf{Input:}}
\algrenewcommand\algorithmicensure{\textbf{Output:}}
\newcommand{\algP}{\ptxt{\textbf{P}}}
\newcommand{\algNP}{\ptxt{\textbf{NP}}}
\newcommand{\algNPC}{\ptxt{\textbf{NP-Complete}}}
\newcommand{\algNPH}{\ptxt{\textbf{NP-Hard}}}
\newcommand{\algEXP}{\ptxt{\textbf{EXP}}}
\DeclareMathOperator{\fl}{fl}

%%%%%%%%%%%%%%%%%%%%%%%%%%%%%
% Other commands
%%%%%%%%%%%%%%%%%%%%%%%%%%%%%
\newcommand{\flag}[1]{\textbf{\textcolor{red}{#1}}}
\let\uSym\u
\let\u\relax
\newcommand{\u}[1]{\underline{#1}}
\let\bSym\b
\let\b\relax
\newcommand{\b}[1]{\textbf{#1}}
\let\iSym\i
\let\i\relax
\newcommand{\i}[1]{\textit{#1}}
\let\scSym\sc
\let\sc\relax
\newcommand{\sc}[1]{\textsc{#1}}

%%%%%%%%%%%%%%%%%%%%%%%%%%%%%%%%%%%%%%%
% Make l's curvy in math environments %
%%%%%%%%%%%%%%%%%%%%%%%%%%%%%%%%%%%%%%%
\mathcode`l="8000
\begingroup
\makeatletter
\lccode`\~=`\l
\DeclareMathSymbol{\lsb@l}{\mathalpha}{letters}{`l}
\lowercase{\gdef~{\ifnum\the\mathgroup=\m@ne \ell \else \lsb@l \fi}}%
\endgroup

%%%%%%%%%%%%%%%%%%%%%%%%%
% Fix \vdots and \ddots %
%%%%%%%%%%%%%%%%%%%%%%%%%
\usepackage{letltxmacro}
\LetLtxMacro\orgvdots\vdots
\LetLtxMacro\orgddots\ddots

\makeatletter
\DeclareRobustCommand\vdots{%
	\mathpalette\@vdots{}%
}
\newcommand*{\@vdots}[2]{%
	% #1: math style
	% #2: unused
	\sbox0{$#1\cdotp\cdotp\cdotp\m@th$}%
	\sbox2{$#1.\m@th$}%
	\vbox{%
		\dimen@=\wd0 %
		\advance\dimen@ -3\ht2 %
		\kern.5\dimen@
		% remove side bearings
		\dimen@=\wd2 %
		\advance\dimen@ -\ht2 %
		\dimen2=\wd0 %
		\advance\dimen2 -\dimen@
		\vbox to \dimen2{%
			\offinterlineskip
			\copy2 \vfill\copy2 \vfill\copy2 %
		}%
	}%
}
\DeclareRobustCommand\ddots{%
	\mathinner{%
		\mathpalette\@ddots{}%
		\mkern\thinmuskip
	}%
}
\newcommand*{\@ddots}[2]{%
	% #1: math style
	% #2: unused
	\sbox0{$#1\cdotp\cdotp\cdotp\m@th$}%
	\sbox2{$#1.\m@th$}%
	\vbox{%
		\dimen@=\wd0 %
		\advance\dimen@ -3\ht2 %
		\kern.5\dimen@
		% remove side bearings
		\dimen@=\wd2 %
		\advance\dimen@ -\ht2 %
		\dimen2=\wd0 %
		\advance\dimen2 -\dimen@
		\vbox to \dimen2{%
			\offinterlineskip
			\hbox{$#1\mathpunct{.}\m@th$}%
			\vfill
			\hbox{$#1\mathpunct{\kern\wd2}\mathpunct{.}\m@th$}%
			\vfill
			\hbox{$#1\mathpunct{\kern\wd2}\mathpunct{\kern\wd2}\mathpunct{.}\m@th$}%
		}%
	}%
}
\makeatother

\newcommand{\B}{
	\begin{tikzpicture}
	\filldraw [fill=red, draw=black] (0, 0) rectangle (0.37, 0.45);
	\draw [line width=0.5mm, white ] (0.1,0.08) -- (0.1,0.38);
	\draw[line width=0.5mm, white ] (0.1, 0.35) .. controls (0.2, 0.35) and (0.4, 0.2625) .. (0.1, 0.225);
	\draw[line width=0.5mm, white ] (0.1, 0.225) .. controls (0.2, 0.225) and (0.4, 0.1625) .. (0.1, 0.1);
	\end{tikzpicture}
}

% Allow custom symbols as arrows in tikz-cd
\tikzset{
  symbol/.style={
    draw=none,
    every to/.append style={
      edge node={node [sloped, allow upside down, auto=false]{$#1$}}}
  }
}
% Example: \arrow[r,symbol=\cong]
% https://tex.stackexchange.com/questions/394154/how-to-include-inclusion-subgroup-relationship-in-tikz-cd-diagram

\author{Professor Alejandro Uribe-Ahumada\\ \small\i{Transcribed by Thomas Cohn}}
\title{Math 591 Lecture 29}
\date{11/9/20} % Can also use \today

\begin{document}
\maketitle
\setlength\RaggedRightParindent{\parindent}
\RaggedRight

\subsection*{Motivation: Integration}

\par\noindent
How do we integrate on a manifold? Start with calc 3. Let $U\subseteq\R^{n}$ be an open set. Then we can change variables with
\[
	\int_{U}f(x)\,dx=\int_{V}f(x(y))\underbrace{\det \begin{pmatrix}
		\pd{x}{y}
	\end{pmatrix}}_{\ptxt{if positive}}\,dy
\]
for $x=x(y)$.\n

\par\noindent
Now, for the general interpretation: let $\map{F:V}{U}{y}{x=x(y)}$ be a diffeomorphism. Then $\begin{pmatrix}\pd{x}{y}\end{pmatrix}$ is the Jacobian of $F$. We'll interpret this as the pullback of the RHS integral by $F$ -- $F^{*}(dx^{1},\ldots,dx^{n})=?$. But what does $\det$ mean in general? It's an alternating, multilinear function, so let's work with that.\n

\defn{
	Let $V$ be a $n$-dimensional vector space, and $k\in\N$. A \u{$k$-covector} or \u{$k$-form} on $V$ is a multilinear map $\alpha:\underbrace{V\times\cdots\times{}V}_{k}\to\R$ that is alternating (i.e. if you swap two elements, the sign flips).\n
}

\par\noindent
Observe: Let $\sigma\in{}S_{k}$ (the symmetric group). Given $v_{1},\ldots,v_{k}$, $\alpha$ alternating, we have
\[
	\alpha(v_{\sigma(1)},\ldots,v_{\sigma(k)})=\underbrace{(-1)^{\sigma}}_{\mathclap{=\sgn(\sigma)}}\alpha(v_{1},\ldots,v_{k})
\]

\par\noindent
Observe: If $\set{v_{1},\ldots,v_{k}}$ is linearly dependent, and $\alpha$ is alternating, then $\alpha(v_{1},\ldots,v_{k})=0$.\n
Proof: Say $v_{1}=\lambda_{2}v_{2}+\cdots+\lambda_{k}v_{k}$. Then
\[
	\alpha(v_{1},\ldots,v_{k})=\alpha\paren{\sum_{i=2}^{k}\lambda_{i}v_{i},v_{2},\ldots,v_{k}}=\sum_{i=2}^{k}\lambda_{i}\alpha(v_{i},v_{2},\ldots,v_{i},\ldots,v_{k})=0
\]

\defn{
	$\bigwedge^{k}V^{*}$ is the set of alternating $\R$-multilinear functions on $V^{k}$. We say that $k$ is the \u{degree}.\n
}

\par\noindent
Observe: $k$-forms can be pulled back by linear maps.\n

\defn{
	If $F:V\to{}W$ is linear, and $\alpha\in{}\bigwedge^{k}W^{*}$, $v_{1},\ldots,v_{k}\in{}V$, then
	\[
		(F^{*}\alpha)(v_{1},\ldots,v_{k})=\alpha(F(v_{1}),\ldots,F(v_{k}))
	\]
	and $F^{*}\alpha\in\bigwedge^{k}V^{*}$.\n
}

\defn{
	Let $(\mathcal{E}^{1},\ldots,\mathcal{E}^{n})$ be an ordered basis of $V^{*}$. Let $A=\set{a_{1}<\cdots<a_{k}}\subset\set{1,\ldots,n}$. Define $\mathcal{E}^{A}:V\times\cdots\times{}V\to\R$ by
	\[
		\mathcal{E}^{A}(v_{1},\ldots,v_{k})=\det \begin{pmatrix}
			\mathcal{E}^{a_{i}}(v_{j})
		\end{pmatrix}_{(i,j)}
	\]
}

\ex{
	Say $V=\R^{3}$ with the standard basis. Then $\mathcal{E}^{13}((x_{1},x_{2},x_{3}),(y_{1},y_{2},y_{3}))=\det \begin{pmatrix}
		x_{1} & y_{1}\\
		x_{3} & y_{3}
	\end{pmatrix}$.\n
}

\prop{
	For a given $V$ and $k$, and an ordered basis $(\mathcal{E}^{1},\ldots,\mathcal{E}^{n})$ of $V^{*}$, the set $\set{\mathcal{E}^{I}:I\subset\set{1,\ldots,n},\#I=k}$ is a basis of $\bigwedge^{k}V^{*}$. In particular, $\dim\bigwedge^{k}V^{*}=\binom{n}{k}$.\n
}

\par\noindent
Observe: If $k>n$, $\bigwedge^{k}V^{*}=\set{0}$. If $k=n$, $\dim\bigwedge^{k}V^{*}=1$. If $k=1$, $\bigwedge^{1}V^{*}=V^{*}$.\n

\newpage
\par\noindent
As a warmup, let $k=2$> Let $\set{e_{1},\ldots,e_{n}}$ be a basis of $V$, and $\set{\mathcal{E}^{1},\ldots,\mathcal{E}^{n}}$ the corresonding dual basis of $V^{*}$. Note that $\mathcal{E}^{i}(e_{j})=\delta_{ij}$.\n
Say $\alpha\in\bigwedge^{2}V^{*}$, $v_{1}=\sum_{a=1}^{n}v_{1}^{a}e_{a}$, $v_{2}=\sum_{b=1}^{n}v_{2}^{b}e_{b}$. Then
\begin{align*}
	\alpha(v_{1},v_{2}) & =\alpha\paren{\sum_{a=1}^{n}v_{1}^{a}e_{a},\sum_{b=1}^{n}v_{2}^{b}e_{b}}\\
	& =\sum_{a=1}^{n}v_{1}^{a}\alpha\paren{e_{a},\sum_{b=1}^{n}v_{2}^{b}e_{b}}\\
	& =\sum_{a=1}^{n}\sum_{b=1}^{n}v_{1}^{a}v_{2}^{b}\alpha(e_{a},e_{b})\\
	& =\smashoperator[r]{\sum_{1\le{}a<b\le{}n}}v_{1}^{a}v_{2}^{b}\alpha(e_{a},e_{b})+v_{1}^{b}v_{2}^{a}\alpha(e_{b},e_{a})\\
	& =\sum_{1\le{}a<b\le{}n}\underbrace{(v_{1}^{a}v_{2}^{b}-v_{1}^{b}v_{2}^{a})}_{\displaystyle\hspace{-6em}\mathrlap{=\left|\begin{matrix}
			v_{1}^{a} & v_{2}^{a}\\ v_{1}^{b} & v_{2}^{b})
		\end{matrix}\right|=\mathcal{E}^{ab}(v_{1},v_{2})}}\alpha(e_{a},e_{b})\\
	& =\sum_{I=\set{a<b}}\alpha(e_{a},e_{b})\mathcal{E}^{I}(v_{1},v_{2})
\end{align*}
In general, $\displaystyle\alpha=\smashoperator{\sum_{1\le{}a<b\le{}n}}\alpha(e_{a},e_{b})\mathcal{E}^{ab}$.\n

\par\noindent
Now, for general $k\in\N$, fix $\alpha\in\bigwedge^{k}V^{*}$. For $i=1,\ldots,k$, let $v_{i}=\sum_{a=1}^{n}v_{i}^{a}e_{a}$. Then
\begin{align*}
	\alpha(v_{1},\ldots,v_{k}) & =\alpha\paren{\sum_{a_{1}=1}^{n}v_{1}^{a_{1}}e_{a_{1}},\ldots,\sum_{a_{k}}^{n}v_{k}^{a_{k}}e_{a_{k}}}\\
	& =\sum_{a_{1},\ldots,a_{k}=1}^{n}\paren{\prod_{j=1}^{k}v_{j}^{a_{j}}}\alpha(e_{a_{1}},\ldots,e_{a_{k}})\\
	& =\sum_{1\le{}a_{1}<\cdots<a_{k}\le{}n}\underbrace{\paren{\sum_{\sigma\in{}S_{k}}\paren{\prod_{j=1}^{k}v_{j}^{\sigma(a_{i})}}(-1)^{\sigma}}}_{\displaystyle\hspace{-12em}\mathrlap{=\det\begin{pmatrix}v_{j}^{a_{i}}\end{pmatrix}_{(i,j)}=\mathcal{E}^{A}(v_{1},\ldots,v_{k}),A=\set{a_{1},\ldots,a_{k}}}}\alpha(e_{a_{1}},\ldots,e_{a_{k}})\\
	& =\smashoperator[r]{\sum_{A=\set{a_{1}<\cdots<a_{k}}}}\alpha(e_{a_{1}},\ldots,e_{a_{k}})\mathcal{E}^{A}(v_{1},\ldots,v_{k})
\end{align*}
So in general, $\displaystyle\alpha=\smashoperator{\sum_{A=\set{a_{1},\cdots<a_{k}}}}\alpha(e_{a_{1}},\ldots,e_{a_{k}})\mathcal{E}^{A}$.\n

\section*{The Wedge Product}

\par\noindent
We want to take a $k$-form and an $l$-form and make a $k+l$-form.\n

\defn{
	For $\alpha\in\bigwedge^{k}V^{*}$, $\beta\in\bigwedge^{l}V^{*}$, define the \u{wedge product} of $\alpha$ and $\beta$, $\alpha\wedge\beta\in\bigwedge^{k+l}V^{*}$, by
	\begin{align*}
		(\alpha\wedge\beta)(v_{1},\ldots,v_{k+l}) & \eqdef\frac{1}{k!l!}\sum_{\sigma\in{}S_{k+l}}(-1)^{\sigma}\alpha(v_{\sigma(1)},\ldots,v_{\sigma(k)})\beta(v_{\sigma(k+1)},\ldots,v_{\sigma(k+l)})\\
		& =\sum_{\sigma\in\Sh(k,l)}(-1)^{\sigma}\alpha(v_{\sigma(1)},\ldots,v_{\sigma(k)})\beta(v_{\sigma(k+1)},\ldots,v_{\sigma(k+l)})\\
	\end{align*}
}

\defn{
	$\Sh(k,l)$ is the set of \u{$k-l$ shuffles}, which are permutations $\sigma\in{}S(k+l)$ such that $\sigma(1)<\cdots<\sigma(k)$ and\break
	$\sigma_{k+1}<\cdots<\sigma_{k+l}$.\n
}

\subsection*{The Skew-Symmetrizer}

\par\noindent
Given $\alpha\in\bigwedge^{k}V^{*}$ and $\beta\in\bigwedge^{l}V^{*}$, we can define
\[
	(\alpha\otimes\beta)(v_{1},\ldots,v_{k+l})=\alpha(v_{1},\ldots,v_{k})\beta(v_{k+1},\ldots,v_{k+l})
\]
As a map, $\alpha\otimes\beta:V^{k+l}\to\R$ is $k+l$-multilinear, but not alternating/skew-symmetric.\n

\defn{
	The \u{skew-symmetrizer} of a multilinear map $f:V^{m}\to\R$ is defined by
	\[
		A(f)(v_{1},\ldots,v_{m})=\frac{1}{m!}\sum_{\sigma\in{}S_{m}}(-1)^{\sigma}f(v_{\sigma(1)},\ldots,v_{\sigma(m)})
	\]
	\n
}

\lemma{
	$A(f)$ is skew-symmetric/alternating, and if $f$ is already skew-symmetric, $A(f)=f$.\nn
	Proof: Let $\tau\in{}S_{m}$. Then
	\begin{align*}
		A(f)(v_{\tau(1)},\ldots,v_{\tau(m)}) & =\frac{1}{m!}\sum_{\sigma\in{}S_{m}}(-1)^{\sigma}f(v_{\sigma(\tau(1))},\ldots,v_{\sigma(\tau(m))})\\
		& \labeledeq{(1)}\frac{1}{m!}\sum_{\mu\in{}S_{m}}\underbrace{(-1)^{\mu\tau}}_{\hspace{-3em}\mathrlap{=(-1)^{\tau}(-1)^{\mu}}}f(v_{\mu(1)},\ldots,v_{\mu(m)})\\
		& =(-1)^{\tau}A(f)
	\end{align*}
	with (1) true because if $\mu=\sigma\tau$, then $\sigma=\mu\tau\inv$.\proven
}

\par\noindent
Thus, we have $\alpha\wedge\beta=\frac{(k+l)!}{k!l!}A(\alpha\otimes\beta)$.\n

\ex{
	$k=l=1$. Then $\Sh(1,1)=S_{2}$, so
	\[
		(\alpha\wedge\beta)(v_{1},v_{2})=\alpha(v_{1})\beta(v_{2})-\alpha(v_{2})\beta(v_{1}).
	\]
}

\ex{
	$k=1$, $l=2$. Then the elements of $S_{3}$ are
	\[
		\begin{array}{ccc|c|c}
			\sigma(1) & \sigma(2) & \sigma(3) & \in\Sh(1,2)? & \sgn\\ \hline
			1 & 2 & 3 & \ptxt{Yes} & +\\
			1 & 3 & 2 & \ptxt{No}\\
			2 & 1 & 3 & \ptxt{Yes} & -\\
			2 & 3 & 1 & \ptxt{No}\\
			3 & 1 & 2 & \ptxt{Yes} & +\\
			3 & 2 & 1 & \ptxt{No}
		\end{array}
	\]
	so $\Sh(1,2)=\set{\cycle{1 & 2 & 3},\cycle{2 & 1 & 3},\cycle{3 & 1 & 2}}$, so
	\[
		(\alpha\wedge\beta)(v_{1},v_{2},v_{3})=\alpha(v_{1})\beta(v_{2},v_{3})-\alpha(v_{2})\beta(v_{1},v_{3})+\alpha(v_{3})\beta(v_{1},v_{2})
	\]
	If $\beta=\gamma\wedge\delta$, then we get
	\begin{align*}
		(\alpha\wedge(\gamma\wedge\delta))(v_{1},v_{2},v_{3})= & \alpha(v_{1})(\gamma(v_{2})\delta(v_{3})-\gamma(v_{3})\delta(v_{2}))\\
		& -\alpha(v_{2})(\gamma(v_{1})\delta(v_{3})-\gamma(v_{3})\delta(v_{1}))\\
		& +\alpha(v_{3})(\gamma(v_{1})\delta(v_{2})-\gamma(v_{2})\delta(v_{1}))\\
		= & \left|\begin{matrix}
			\alpha(v_{1}) & \alpha(v_{2}) & \alpha(v_{3})\\
			\gamma(v_{1}) & \gamma(v_{2}) & \gamma(v_{3})\\
			\delta(v_{1}) & \delta(v_{2}) & \delta(v_{3})
		\end{matrix}\right|
	\end{align*}
}

\lemma{
	$\mathcal{E}^{I}\wedge\mathcal{E}^{J}=\mathcal{E}^{IJ}$.\n
}

\lemma{
	$\mathcal{E}^{I}=\mathcal{E}^{I_{1}}\wedge\cdots\wedge\mathcal{E}^{I_{k}}$ ($\wedge$ is associative).\n
}

\prop{
	The wedge product is
	\begin{itemize}[leftmargin=4\parindent]
		\item Bilinear
		\item Associative
		\item Anti-commutative: $\beta\wedge\alpha=(-1)^{kl}\alpha\wedge\beta$, for $\alpha\in\bigwedge^{k}$, $\beta\in\bigwedge^{l}$
		\item If $F:V\to{}W$, $\alpha\in\bigwedge^{k}W^{*}$, $\beta\in\bigwedge^{l}W^{*}$, then $F^{*}(\alpha\wedge\beta)=(F^{\alpha})\wedge(F^{*}\beta)$
	\end{itemize}
}

\end{document}